\section{Probleme}

\subsection{Problema Cut and Paste (CSES)}
\label{problem:cut-and-paste}

\href{https://cses.fi/problemset/task/2072}{enunț}
$\bullet$
\hyperref[code:cut-and-paste]{sursă}

Problema este o aplicație directă a treapurilor implicite. Creăm un treap cu o literă în fiecare nod. Pentru o operație $(a, b)$, este suficient să spargem treapul în trei porțiuni, respectiv $X = [1, a-1]$, $Y = [a, b]$ și $Z = [b + 1, n]$. Apoi alipim aceste porțiuni în ordinea $XZY$.

Remarcăm că această problemă se mulează perfect peste implementarea cu \textit{split} și \textit{merge}, dar nu și peste cea cu rotații.

\subsection{Problema Reversals and Sums (CSES)}
\label{problem:reversals-and-sums}

\href{https://cses.fi/problemset/task/2074}{enunț}
$\bullet$
\hyperref[code:reversals-and-sums]{sursă}

Această problemă este un exemplu tipic de treap cu informații cu propagare \textit{lazy}. Merită să luăm în calcul posibilitatea de a folosi un simplu arbore de intervale, dar din păcate acela nu se pretează. Dacă descompunem un interval $[a,b]$ și notăm în fiecare nod din descompunere informația \textit{lazy} booleană „acest interval trebuie răsturnat”, vom greși. Intervalele nu trebuie doar răsturnate, ci și mutate la stînga sau la dreapta, căci vor ajunge să ocupe alte poziții în listă.

Să pornim de la un treap implicit, căci conceptual memorăm valori dintr-un vector. Acest treap vine cu cîmpurile obișnuite în fiecare nod: valoare, prioritate și contor (mărimea subarborelui). În plus, vom adăuga încă două cantități:

\begin{enumerate}
  \item O valoare întreagă \ccode{subtree_sum} care reține suma valorilor din subarbore. Vom menține în timp real această sumă, ca și pe contor, într-o funcție \ccode{update_node} pe care o apelăm ori de cîte ori structura arborelui se modifică. Știm deja să facem acest lucru pentru treapuri implicite.

  \item O valoare booleană \ccode{flip} care indică că întregul subarbore trebuie răsturnat.
\end{enumerate}

Să vedem acum cum implementăm cele două operații.

Pentru \ccode{range_sum(a, b)}, putem să spargem treapul ca și în problema anterioară în $X = [1, a-1]$, $Y = [a, b]$ și $Z = [b + 1, n]$. Răspunsul va fi valoarea \ccode{subtree_sum} din $Y$. Apoi alipim la loc fragmentele în ordinea inițială.

Pentru \ccode{range_reverse(a, b)}, procedăm la fel, dar pe rădăcina lui $Y$ setăm \ccode{flip} la valoarea 1.

Ca și la arborii de intervale, este important să propagăm informația \textit{lazy} la fii în momentele-cheie. Am indicat aceste momente prin comentarii în cod. Am denumit funcția tot \ccode{push()} ca și la arborii de intervale.

\subsection{Problema T-Shirts (Codeforces)}
\label{problem:t-shirts}

\href{https://codeforces.com/contest/702/problem/F}{enunț}
$\bullet$
\hyperref[code:t-shirts]{sursă}

Chiar dacă găsim problema în vreo colecție de linkuri pe Internet și știm că se rezolvă cu treapuri, folosirea acestora nu este evidentă.

Să începem cu partea simplă. Pare natural să ordonăm oferta de tricouri așa cum le evaluează clienții: descrescător după calitate, apoi crescător după preț. De aici încolo, calitatea nu ne mai interesează. De asemenea, dacă ne ajută putem ordona bugetele clienților, să spunem crescător.

Am încercat să construiesc un treap implicit cu costurile tricourilor, iar în acest treap, dat fiind un buget $b$, să navighez logaritmic spre ultimul tricou cumpărat. Dar problema este că prețurile nu sînt ordonate, iar clienții nu cumpără tricouri dintr-un interval compact, ci pe sărite. Ne-am dori să aflăm rapid cîte tricouri cumpără un client, și cu ce preț total, din fiul stîng al nodului curent (ca să știm cu ce buget rămas coborîm în fiul drept). Dar pare imposibil. Pentru prețuri ca $1, \np{1000000}, 1, \np{1000000}, 1, \np{1000000}$ nu putem evalua ușor situația.

Problema nu include un editorial, dar am găsit ideea rezolvării în \href{https://codeforces.com/blog/entry/46324#comment-307625}{acest comentariu}. Ținem \textbf{un treap de bugete}, inițializat cu bugetele inițiale ale clienților, ordonate. Apoi evaluăm tricourile în ordinea de mai sus. Întrebarea este cum actualizăm rapid clienții care mai au bani să cumpere tricoul $i$. Dorim să adăugăm $+1$ la numărul de tricouri și $-c_i$ la bugetul rămas pentru toți acești clienți.

Pentru tricoul curent $i$, de cost $c_i$, să apelăm \ccode{split(c[i])}. Am numit treapurile rezultate $T_N$ (de la \textit{no buy} / nu cumpăra) și respectiv $T_B$ (de la \textit{buy} / cumpără). Pe $T_B$ notăm informația \textit{lazy} că toți clienții au $+1$ tricou și $-c_i$ la buget. Dar acum trebuie să refacem treapul ordonat după noile bugete. Cum să procedăm? Operația \ccode{merge()} funcționează doar cînd $T_N$ și $T_B$ acoperă intervale disjuncte de chei (orice cheie din $T_N$ este mai mică sau egală cu orice cheie din $T_B$). Dar, după ce toți clienții din $T_B$ au plătit $c_i$, cheile pot fi intercalate.

Soluția este să spargem din nou $T_B$ în valori $<c_i$ și $\geq c_i$ (pe noile bugete). Să numim aceste treapuri $T_P$ (de la \textit{poor} / sărac) și $T_R$ (de la \textit{rich} / bogat). În cuvinte,

\begin{itemize}
  \item $T_N$ sînt clienți cu bugete din $[0,c_i)$, care nu au cumpărat tricoul $i$.
  \item $T_P$ sînt clienți cu bugete din $[0,c_i)$ după ce au cumpărat tricoul $i$.
  \item $T_R$ sînt clienți cu bugete din $[c_i,\infty)$ după ce au cumpărat tricoul $i$.
\end{itemize}

Acum unificăm $T_N$, $T_P$ și $T_R$ astfel:

\begin{itemize}
  \item Inserăm naiv valorile din $T_P$ în $T_N$. Obținem un treap cu toate valorile în $[0, c)$.
  \item Unificăm acest treap cu $T_R$ apelînd \ccode{merge()}, căci acum cheile sînt separate de $c$.
\end{itemize}

Pasul  inserării naive îl rezolvăm cu un DFS prin $T_P$. Nu este nevoie să implementăm ștergeri din $T_P$, căci îl dezmembrăm complet. Este suficient să setăm pe \ccode{NULL} pointerii fiecărui nod din $T_P$ pentru a-l transforma într-o frunză.

Toate operațiile sînt $\bigoh(n \log n)$, cu excepția inserărilor naive. Care este efortul total pentru aceste operații? Să considerăm ce se întîmplă într-o inserare naivă. Avem un cost $c$ și un buget $b \in [c, 2c)$ (dacă am fi avut $b \geq 2c$, atunci nodul ar fi stat în $T_R$, nu în $T_P$). Dar atunci să observăm că noul buget, $b-c$, scade sub jumătate din bugetul original $b$:

$$b < 2c \implies \frac{b}{2} < c \implies b - c < \frac{b}{2}$$

În cuvinte, de cîte ori este inserat naiv în $T_N$, bugetul clientului scade la jumătate sau chiar mai jos. Deci fiecare client poate fi inserat naiv de cel mult $\log VAL\_MAX$ ori. Complexitatea inserărilor naive, și a întregului algoritm, este $\bigoh(n \log n \log VAL\_MAX)$.

\subsubsection*{Detaliu de implementare}

În cod veți observa că funcțiile \ccode{split()} și \ccode{insert()} folosesc operatorii $<$ și respectiv $\leq$ pentru a compara bugetele. Această decizie este importantă, iar \href{https://codeforces.com/blog/entry/92340#comment-814503}{acest comentariu} de pe Codeforces explică foarte bine motivul. Pe măsură ce procesăm tricouri, iar bugetele clienților încep să scadă, vom avea tot mai multe bugete egale. În acest moment trebuie să fim atenți cum gestionăm cheile egale într-un treap. Altfel riscăm să dezechilibrăm treapul și să ajungem la adîncime $\bigoh(n)$.

Să ne imaginăm o operație \textit{split($x$)} pe un treap cu toate cheile egale cu $x$. Să spunem că algoritmul alege să mute cheile egale cu $x$ în dreapta, deci după spargere treapul stîng va rămîne gol.

Acum să ne imaginăm operația \textit{insert($x$)}, din nou într-un treap cu toate cheile egale cu $x$. După cum știm, \textit{insert} începe prin a coborî de la rădăcină pînă cînd prioritatea noului nod depășește prioritatea subarborelui curent. Toate cheile fiind egale, noi decidem dacă lăsăm funcția să coboare în nodul stîng sau în cel drept. Cînd prioritatea noului nod este maximă, spargem subarborele în două treapuri care devin fiii nodului inserat.

Aici intervine dezechilibrarea. Dacă lăsăm funcția \textit{insert} să coboare în fiul drept, iar dacă \textit{split} transferă toate nodurile existente tot în fiul drept, putem obține în timp un treap de adîncime $\bigoh(n)$. Este important ca \textit{insert} să prefere direcția opusă lui \textit{split}.


\subsection{Problema Clasament (Lot 2016)}
\label{problem:clasament}

\href{https://kilonova.ro/problems/1972}{enunț}
$\bullet$
\hyperref[code:clasament]{surse}

Cred că problema „se cere” rezolvată cu hashing. Nu știu cum altfel putem compara și clasifica $m$ permutări de lungime $n$. Dacă acceptăm asta, atunci cumva trebuie să menținem eficient codul hash al permutării. Un algoritm de hashing eficient și ușor de implementat este să privim permutarea ca pe un număr de $n$ cifre într-o bază $B$ și să calculăm acel număr modulo un $m$ ales.

\subsubsection*{Soluție cu treap}

Fiindcă sîntem la capitolul de treapuri, să ținem într-un treap ID-urile elevilor (de la $0$ la $n-1$), ordonate după scor și timp. Atunci putem calcula recursiv codul hash al unui treap. Dacă treapul stochează ID-urile 5, 8 și 2 în subarborele stîng, 4 în rădăcină și 3 și 7 în subarborele drept, și dacă luăm $B=10$ pentru simplitate, atunci dorim ca codul hash să fie 582437, pe care îl putem sparge ca:

$$582437 = 582 \cdot 10^3 + 4 \cdot 10^2 + 37$$

Generalizînd,

\begin{itemize}
  \item Fie $L$ codul hash al subarborelui fiului stîng.
  \item Fie $R$ codul hash al subarborelui fiului drept.
  \item Fie $S_R$ mărimea (numărul de noduri) al subarborelui fiului drept.
  \item Fie $i$ ID-ul din rădăcina subarborelui.
  \item Atunci codul hash este $(L \cdot B^{S_R + 1} + i \cdot B^{S_R} + R) \bmod M$
\end{itemize}

Este important să alegem o bază $B$ primă. Altfel putem ajunge la coliziuni (permutări diferite cu coduri hash identice).

Acum putem definitiva structura de date. În fiecare nod avem nevoie să stocăm codul hash și mărimea subarborelui. Putem recalcula aceste valori în $\bigoh(1)$ cînd reorganizăm treapul, fie că alegem implementarea cu \textit{split} și \textit{merge}, fie cu rotații.

Ca detaliu de implementare, cînd un scor se modifică trebuie să ștergem ID-ul elevului respectiv din treap și să-l reinserăm conform noului scor. Ca să nu risipim memorie, cu puțină atenție putem refolosi chiar nodul șters.

\subsubsection*{Soluție cu descompunere în radical}

Putem să stocăm permutarea curentă și în blocuri de mărime $\bigoh(\sqrt{n})$. Această descompunere amintește de problema \hyperref[problem:serega-and-fun]{Serega and Fun} de la capitolul de descompunere în radical. Acolo aveam nevoie să facem rotații pe subsecvențe, dar operația din problema de față este similară.

Nu insistăm ca blocurile să aibă fix lungime $\sqrt{n}$, ci le lăsăm să se dezechilibreze prin ștergeri și inserări. Dacă un bloc ajunge la o lungime-limită, să spunem $2 \sqrt{n}$, atunci reechilibrăm structura: colectăm toate blocurile într-un vector, apoi le redistribuim. În practică, testele problemei nu duc la reechilibrări.

Pe fiecare bloc menținem codul hash al fragmentului de permutare din acel bloc. La o modificare de scor, recalculăm în $\bigoh(\sqrt{n})$ codul hash al celor (cel mult) două blocuri afectate. După fiecare modificare, recalculăm codul hash al întregii permutări combinînd codurile celor $\bigoh(\sqrt{n})$ fragmente după aceeași regulă ca mai sus.

Cele două soluții au timpi practic egali, ceea ce arată că treapurile au constantă cam proastă.


\subsection{Problema Roata Norocului (Lot 2024)}
\label{problem:roata-norocului}

\href{https://kilonova.ro/problems/2804}{enunț}
$\bullet$
\hyperref[code:roata-norocului]{sursă}

\subsubsection*{Costul comprimării unui arbore}

Să considerăm cazul unui arbore-stea cu 5 noduri, cu centrul de valoare 8 și vecinii de valori 6, 10, 20 și 30. Să analizăm două strategii. În prima, unim 8 cu vecinii mai scumpi, apoi cu cel mai ieftin.

\begin{itemize}
  \item Unește 10-8, cost 2.
  \item Unește 20-8, cost 12.
  \item Unește 30-8, cost 22.
  \item Unește 8-6, cost 2.
  \item Cost total 38.
\end{itemize}

În a doua strategie, unim 8 cu 6 cîndva mai devreme:

\begin{itemize}
  \item Unește 10-8, cost 2.
  \item Unește 8-6, cost 2.
  \item Unește 20-6, cost 14.
  \item Unește 30-6, cost 24.
  \item Cost total 42.
\end{itemize}

Pare că rentează să unim fiecare nod cu vecinii săi mai scumpi decît el înainte de a-l uni cu unul mai mic. Pentru a ne verifica intuiția, să privim coloana stîngă în listele de mai sus. Toate valorile apar exact o dată, cu excepția minimului (6). Este normal, căci un nod apare în stînga în momentul cînd este eliminat. Așadar, suma coloanei stîngi este constantă. Atunci, pentru a minimiza costul, trebuie să maximizăm coloana dreaptă. De aceea vom amîna să eliminăm un nod pînă cînd devine inevitabil.

(Ca fapt divers, problema în forma inițială se referea la un singur arbore, fără actualizări. În acel caz, o strategie suficientă este: la fiecare pas, comprimăm muchia care maximizează minimul dintre cele două noduri.)

Rezultă că fiecare nod va contribui la coloana dreaptă de atîtea ori cîți vecini mai scumpi decît el are. La egalitate putem departaja oricum, cu condiția să nu facem dublă contabilitate, ci exact unul dintre noduri să-l contorizeze pe celălalt. Deci formula pentru cost este:

$$cost(u) = \sum_{x\ \text{în subarbore}} V_x - \min_{x\ \text{în subarbore}} V_x - \sum_{x\ \text{în subarbore}}V_x \cdot B_x$$

Am notat cu $B_x$ numărul de vecini ai lui $x$ mai scumpi decît $x$. În funcție de implementare, poate fi nevoie să scădem 1 din $B_x$ însuși, pentru a-l exclude pe părintele lui $x$. Implementarea de mai jos evită acest caz.

\subsubsection*{Operațiile necesare}

Oricînd dorim să proiectăm o structură de date, pornim de la întrebarea: ce vrem să aflăm?

Primii doi termeni ai formulei, adică suma și minimul din subarbore, cu actualizări, sînt partea ușoară a problemei. Vom face o \hyperref[sec:linearization-types]{liniarizare DFS}, peste care construim un arbore de intervale de sume și de minime, cu actualizări punctuale și cu interogări pe interval. Dacă notăm cu $t_i[u]$ și $t_o[u]$ timpii de intrare și de ieșire din DFS, atunci aflăm primii doi termeni ai formulei făcînd interogări pe intervalul $[t_i[u], t_o[u]]$.

Pentru ultimul termen al formulei, eu am reorganizat informația astfel. Fiecare muchie $(x,y)$ are o \textbf{contribuție}, desigur egală cu $\min(V_x, V_y)$. Dacă însumăm contribuțiile tuturor muchiilor din subarbore, ajungem exact la suma dorită: valoarea $V_x$ luată de atîtea ori de cîte ori $x$ este minim pe o muchie.

Fiecare nod $u$ este responsabil să stocheze \textbf{suma contribuțiilor pe muchiile de la $u$ la fiii lui $u$}. Să numim această cantitate \textbf{contribuția unui nod}. Putem stoca contribuțiile nodurilor tot în arborele de intervale. Atunci contribuția totală a unui subarbore (acel ultim termen al formulei) este tot o simplă sumă pe interval.

\subsubsection*{Interogări}

Răspunul la o interogare este așadar cumulul a trei valori într-un arbore de intervare elementar: suma valorilor, minus valoarea minimă, minus suma contribuțiilor.

\subsubsection*{Actualizări}

Desigur, actualizarea valorii unui nod în arborele de intervale este elementară. Dar cum actualizăm contribuția unui nod? Să vedem ce contribuții se schimbă pe muchii la modificarea valorii unui nod $u$ din $V_u$ în $V'_u$.

\begin{enumerate}
  \item Suma contribuțiilor dintre $u$ și fiii lui $u$, stocată în $u$.
  \item Contribuția dintre $u$ și părintele său, stocată în părinte (dacă există).
\end{enumerate}

Pentru (1), să presupunem că valoarea crește, deci $V_u < V'_u$. Cazul în care valoarea scade se calculează la fel, doar cu semn schimbat.

\begin{itemize}
  \item Pentru fiii $x$ cu $V_x < V_u$ contribuția nu se schimbă: contribuția pe acele muchii continuă să fie $V_x$.

  \item Pentru fiii $x$ cu $V_u \leq V_x < V'_u$ contribuția era $V_u$ și acum devine $V_x$. Să spunem că există $k$ astfel de fii. Pentru a actualiza contribuția nodului $u$, vom scădea $k \cdot V_u$ și vom adăuga suma valorilor acestor fiii.

  \item Pentru fiii $x$ cu $V'_u < V_x$ contribuția era $V_u$ și acum devine $V'_u$. Să spunem că există $l$ astfel de fii. Pentru a actualiza contribuția nodului $u$, vom aduna $l \cdot (V'_u - V_u)$.
\end{itemize}

Așadar, avem nevoie de o structură care să ne dea rapid numărul de elemente dintr-un interval de valori dat și suma valorilor acestora.

Pentru (2), putem calcula naiv contribuția, fiind vorba de o singură muchie. Apoi este nevoie să-i spunem părintelui: în lista ta de valori se află o valoare $V_u$. Schimbă acea valoarea în $V'_u$.

Din păcate, nu cunosc nicio structură de date pentru aceste nevoie mai simplă decît un treap. Pe lîngă augmentarea nodurilor din treap cu numărul de elemente din subarbore, avem nevoie și de augmentarea cu suma elementelor. Întrucît fiecare nod își construiește un treap cu fiii, suma mărimilor treapurilor este $n - 1$.
